% =====================================================================
%  Chapter: Unified Evaluation
% =====================================================================
\chapter{Unified Evaluation}
\label{ch:evaluation}

The value of this thesis lies less in any single method than in their
comparison \emph{on equal terms}: the same SPX days, the same hold-out
protocol, the same arbitrage audit. This chapter assembles the per-day
outputs of the preceding chapters
(\texttt{benchmark\_svi\_slices\_\{full,paper\}.parquet},
\texttt{benchmark\_ssvi\_days\_\{full,paper\}.parquet},
\texttt{deep\_smoother\_days.parquet}, \texttt{operator\_days.parquet}) into
a single evaluation.

\section{Protocol}
Five surfaces are compared on every evaluation day: per-slice \textbf{SVI},
surface \textbf{SSVI}, the proprietary \textbf{OptionMetrics} smoothed
surface, the per-day \textbf{deep smoother} of \cref{ch:deepsmoother}, and
the \textbf{neural operator} of \cref{ch:operator} (evaluated only on its
chronological test window, where it has never seen the day). Metrics:
median IV RMSE on held-out quotes (vol points); butterfly-violation rate and
minimum Durrleman $g$ on a dense audit grid; cross-maturity crossedness;
and the computational cost of producing a new surface.

\section{Headline comparison}

\begin{table}[t]
  \centering
  \caption[Unified five-method comparison]{Unified comparison (median over
  evaluation days; hold-out IV RMSE in vol points). Preliminary values from
  the 8-day sample where available; \TODO{fill remaining cells from the
  full-sample runs of NB02--NB05 before submission.}}
  \label{tab:unified}
  \begin{tabular}{lS[table-format=1.3]S[table-format=2.1]l}
    \toprule
    Method & {Hold-out RMSE} & {Bfly viol.\ (\%)} & Cost of a new surface \\
    \midrule
    SVI (per slice)        & 0.968 & 6.2 & one calibration per slice \\
    SSVI (per day)         & 2.790 & 0.0 & one 3-parameter fit per day \\
    OptionMetrics surface  & {---} & {\TODO{}} & proprietary \\
    Deep smoother (per day)& {\TODO{}} & {\TODO{}} & one training run per day \\
    Neural operator        & {\TODO{}} & {\TODO{}} & one forward pass \\
    \bottomrule
  \end{tabular}
\end{table}

Interpretation of the headline table follows the arc of the thesis. The gap
between SVI and SSVI ($\approx0.97$ vs $\approx2.8$ vol points hold-out) is
the price of built-in no-arbitrage in the parametric world. The deep
smoother is designed to close most of that gap while keeping violations near
zero; the operator is designed to do so \emph{at negligible marginal cost per
day}. Against the industrial reference, our transparent SVI pipeline already
sits within $0.93$ vol points (median) of the OptionMetrics surface over
$2{,}958$ grid points.

\section{Analysis by market regime}
The 2018--2022 sample spans three qualitatively different regimes: the calm
of 2019, the COVID-19 crash of March 2020, and the 2022 rate-hike sell-off.
All per-day metrics are re-aggregated within each regime.
\TODO{insert regime table and figure from NB05 once the full-sample runs
complete; the working hypothesis, supported by the bucket analysis of
\cref{tab:buckets}, is that method rankings are stable but absolute errors
and violation rates degrade sharply in March 2020, and that the operator's
test window (2022) constitutes a genuine out-of-regime test.}

\section{Robustness to sparsity}
The synthetic sparsity study of \cref{sec:dsablation} is repeated on real
days by subsampling quotes before evaluation. Per-slice SVI degrades and
eventually loses calibrable slices; the surface-level methods degrade
gracefully; the operator additionally demonstrates stability of its output
under input subsampling (\cref{sec:invariance}) --- a property no per-surface
method possesses.

